\chapter{Recommandations et Durcissement}
\label{chap:recommandations}

\section{Durcissement administrateur}

\subsection{Trusted Hosts}

L'accès à l'interface d'administration doit être restreint aux seules adresses autorisées :

\begin{lstlisting}[language=bash, caption={Configuration Trusted Hosts}]
config system global
    set admin-trust-hosts enable
    set admin-host "192.168.10.0/24"
end
\end{lstlisting}

\subsection{Protocoles désactivés}

Les protocoles non sécurisés (HTTP, Telnet) doivent être désactivés :

\begin{lstlisting}[language=bash, caption={Désactivation des protocoles non sécurisés}]
config system global
    set admin-http disable
    set admin-telnet disable
    set admin-ssh-port 22
end
\end{lstlisting}

\subsection{Politique de mots de passe}

\begin{lstlisting}[language=bash, caption={Politique de mots de passe}]
config system global
    set password-expir-warning-days 14
    set password-min-len 8
    set password-requires-upper-lower-number-symbol enable
end
\end{lstlisting}

\section{RBAC (Role-Based Access Control)}

\subsection{Profils administrateur}

FortiGate propose des profils prédéfinis pour gérer les droits d'accès :

\begin{table}[H]
\centering
\caption{Profils administrateur recommandés}
\label{tab:rbac}
\begin{tabularx}{\textwidth}{|l|X|X|}
\hline
\textbf{Profil} & \textbf{Droits} & \textbf{Utilisation} \\
\hline
super\_admin\_prof\_admin & Accès total & Administrateur réseau \\
\hline
prof\_admin\_read\_only & Lecture seule & Audit, supervision \\
\hline
prof\_admin\_fw\_only & Pare-feu uniquement & Opérateur sécurité \\
\hline
\end{tabularx}
\end{table}

\subsection{Moindre privilège}

Chaque utilisateur ne doit disposer que des droits strictement nécessaires à ses fonctions. La création d'un profil personnalisé est recommandée :

\begin{lstlisting}[language=bash, caption={Création d'un profil personnalisé}]
config system accprofile
    edit "lab-operator"
        set secfabgrp read-write
        set ftgd-wf read-write
        set vpngrp read-write
        set loggrp read-write
        set fwgrp read-only
        set utmgrp read-write
        set netgrp read-only
    next
end
\end{lstlisting}

\section{Durcissement des interfaces}

\subsection{Services désactivés}

Sur chaque interface, seuls les services nécessaires doivent être activés :

\begin{lstlisting}[language=bash, caption={Durcissement interface WAN}]
config system interface
    edit "port1"
        set allowaccess ping  # Uniquement ping depuis WAN
        set dedicated-to-management no
    next
end
\end{lstlisting}

\subsection{Mode explicit allow}

La politique par défaut doit être « deny all » avec des règles explicites :

\begin{lstlisting}[language=bash, caption={Politique deny explicite}]
config firewall policy
    edit 0
        set name "deny-all"
        set srcintf "any"
        set dstintf "any"
        set action deny
        set schedule "always"
        set service "ALL"
        set logtraffic all
    next
end
\end{lstlisting}

\section{Conformité des logs}

\subsection{Rétention des logs}

Les logs doivent être conservés pour une durée minimale conforme aux exigences :

\begin{lstlisting}[language=bash, caption={Configuration de la rétention des logs}]
config log syslogd setting
    set status enable
    set server "192.168.10.50"
    set port 514
    set mode UDP
end
\end{lstlisting}

\subsection{Alertes}

La configuration d'alertes permet de détecter les événements critiques :

\begin{lstlisting}[language=bash, caption={Configuration des alertes}]
config alertemail setting
    set status enable
    set mail-server "smtp.axeli.ma"
    set mail-prefix "[FortiGate Alert]"
    set severity-filter critical
end
\end{lstlisting}

\section{Recommandations complémentaires}

\begin{enumerate}
    \item \textbf{Mises à jour régulières} : maintenir FortiOS à jour (dans les limites de la licence)
    \item \textbf{Sauvegardes} : export de la configuration avant toute modification majeure
    \item \textbf{Audit régulier} : vérification périodique des politiques et des logs
    \item \textbf{Formation} : maintien des compétences de l'équipe sécurité
    \item \textbf{Plan de réponse} : procédures en cas d'incident de sécurité
\end{enumerate}
